\section{归一化流用可逆性换取精确似然}
\label{sec:14-Deep-Learning/07-Generative-Models-Unified/02}

你要给产线上的视觉异常检测换一个判据。现在用的是 VAE 的重构误差，阈值是拍出来的，换一批相机就得重调；有人提议改成判别器打分，可打分的量纲跟着对手网络走，昨天的 \(0.7\) 和今天的 \(0.7\) 不是一回事。你想要的其实很朴素：一个真正的 \(\log p_\theta(x)\)。它有确定的单位（比特或奈特），可以在不同模型之间比较，可以拿去做似然比，可以直接写进报告。

VAE 给不了——它只给下界，而下界与真值的间隙未知，两个模型的 ELBO 谁高谁低不能直接推出似然谁高谁低。GAN 更彻底，它压根没有密度这个对象。前一篇说过，可算的似然、可采样的模型、任意灵活的架构三样拿不全。归一化流的选择是死守前两样，把第三样交出去。

\subsection{密度不是随便搬运的}

想法很简单：取一个什么都会算的基分布，比如 \(z\sim\mathcal N(0,I_d)\)，再拿一个网络 \(g_\theta\) 把它搬到数据空间，\(x=g_\theta(z)\)。采样立刻就有了——抽 \(z\)，过一遍网络。问题是搬过去之后的密度是多少。

一维上先看清楚。设 \(Z\) 在 \([0,1]\) 上均匀，令 \(Y=(Z+Z^2)/2\)，这是 \([0,1]\) 到 \([0,1]\) 的单调映射。原来 \([0,0.25]\) 这一段装着 \(0.25\) 的概率质量，映射之后变成 \([0,0.156]\)，还是那 \(0.25\) 的质量，但装在更窄的区间里——密度必须升上去。区间被压缩到原来的 \(0.625\) 倍，密度就升到 \(1/0.625=1.6\) 倍。末尾那一段 \([0.75,1]\) 被拉伸到 \([0.656,1]\)，宽了 \(1.375\) 倍，密度相应降到 \(0.727\)。

\begin{center}
\begin{tikzpicture}[x=3.6cm,y=3.6cm,>=stealth]
  \draw[->] (-0.02,0) -- (1.14,0) node[right] {\(z\)};
  \draw[->] (0,-0.02) -- (0,1.16) node[above right] {\(y=f(z)\)};
  \draw[black,thick] plot[domain=0:1,samples=60] (\x,{(\x+\x*\x)/2});
  \foreach \a/\b in {0.25/0.156, 0.5/0.375, 0.75/0.656, 1/1}{
    \draw[gray!55,dashed] (\a,0) -- (\a,\b) -- (0,\b);
  }
  % 输入密度：四段等宽等高
  \foreach \a/\b in {0/0.25, 0.25/0.5, 0.5/0.75, 0.75/1}{
    \draw[fill=blue!16,draw=blue!60!black] (\a,-0.09) rectangle (\b,-0.41);
  }
  % 输出密度：宽度不等，高度反向变化，面积不变
  \foreach \l/\h/\w in {0/0.156/0.512, 0.156/0.375/0.365, 0.375/0.656/0.284, 0.656/1/0.233}{
    \draw[fill=red!14,draw=red!65!black] (-0.09,\l) rectangle ({-0.09-\w},\h);
  }
  \node[blue!60!black,font=\small] at (0.5,-0.53) {输入密度：四段等宽等高};
  \node[red!65!black,font=\small,align=center] at (-0.55,1.20) {输出密度\\宽窄换了，高矮反向换};
  \node[font=\small,align=left,fill=white,inner sep=2pt] at (0.82,0.24) {斜率大处\\密度被摊薄};
  \draw[gray!70] (0.86,0.36) -- (0.92,0.82);
\end{tikzpicture}

\emph{同一块概率质量换了容器：容器变窄多少倍，密度就变高多少倍。Jacobian 行列式记的就是这笔账，它是局部体积的变化率。}
\end{center}

多维上，``区间宽度''换成``体积''，``斜率''换成 Jacobian 行列式。设 \(f_\theta=g_\theta^{-1}\) 是从数据空间回到基空间的映射，则

\[
p_\theta(x)=p_Z\!\left(f_\theta(x)\right)\cdot\abs{\det J_{f_\theta}(x)},
\qquad
J_{f_\theta}(x)=\frac{\partial f_\theta}{\partial x}(x).
\]

取对数：

\[
\log p_\theta(x)=\log p_Z\!\left(f_\theta(x)\right)+\log\abs{\det J_{f_\theta}(x)}.
\]

右边两项都是闭式的：第一项是标准正态的对数密度，一行代码；第二项要求能算 Jacobian 行列式。没有下界，没有采样估计，没有未知间隙——这就是那个真正的 \(\log p_\theta(x)\)。训练就是直接对它做梯度上升，前一篇的谱系里，这是站在前向 KL 方向上、间隙为零的那一条。

\subsection{条件对账：哪些是可以放松的，哪些不是}

变量替换公式要求 \(f_\theta\) 是双射，且 \(f_\theta\) 与其逆都几乎处处可微。这三条各自挡住一种失效。

双射挡住的是质量重复计数。若两个不同的 \(x\) 被映到同一个 \(z\)，那一点的基密度就要被两处共享，公式里的单项表达式直接失效。可微挡住的是体积变化率不存在——行列式是线性化之后的量，没有线性化就没有这个数。而``几乎处处''这个限定说明可以有零测的坏点集，ReLU 那样的分段线性可逆映射并不被排除。

维数不变这一条最容易被当成技术细节想去绕开，但它绕不开。若 \(f_\theta:\R^d\to\R^k\) 而 \(k<d\)，\(J\) 不是方阵，行列式根本没有定义；若 \(k>d\)，\(g_\theta\) 的像集是 \(\R^k\) 里的一张 \(d\) 维流形，Lebesgue 测度为零，推前分布相对 Lebesgue 测度没有密度——你能算的所谓``密度''不存在，不是难算，是不存在。这里的判据是分布之间的绝对连续性，\hyperref[sec:12-Real-Analysis-Support/05-Measure-Integration-and-Conditional-Expectation/04]{``Radon--Nikodym 导数把相对于谁说清楚''} 把它讲清楚了。所以流不能降维：潜空间必须和数据空间一样大，\(32\times32\times3\) 的图像就要 \(3072\) 维的基分布，一个像素的信息都不能压掉。

还有一条不写在公式里但同样跑不掉的限制：连续可逆映射保持连通性。基分布是单峰连通的高斯，它的像也就必然连通。如果真实数据分布的支撑分成几团互不相连的块（不同类别的图像很像这种情形），流只能在块与块之间铺一层薄薄的密度把它们连起来，靠极大的局部拉伸把这层薄膜压到接近零。堆更多层不改变这个事实，只让薄膜更薄、Lipschitz 常数更大、训练更不稳。这是架构层面的表达力上限，不是优化问题。

\subsection{行列式必须便宜}

第二笔代价来自计算。一般 \(d\times d\) 矩阵的行列式要 \(O(d^3)\)。\(d=3072\) 时单次就是约 \(2.9\times10^{10}\) 次运算，而这是\emph{每个样本每一层}都要付的。直接算等于放弃。

耦合层的解法是把矩阵变成三角阵。把输入切成两半 \(x=(x_a,x_b)\)，前一半原样通过，后一半用前一半算出的参数做仿射变换：

\[
\begin{aligned}
y_a&=x_a,\\
y_b&=x_b\odot\exp\!\left(s_\theta(x_a)\right)+t_\theta(x_a).
\end{aligned}
\]

\(s_\theta\) 和 \(t_\theta\) 可以是任意复杂的网络，深度、宽度、注意力，随便。Jacobian 按 \((a,b)\) 分块写出来是

\[
J=\begin{pmatrix} I & 0\\[2pt] \dfrac{\partial y_b}{\partial x_a} & \diag\!\left(\exp s_\theta(x_a)\right)\end{pmatrix},
\]

分块下三角，行列式只由对角块决定，左下那个复杂的块完全不参与：

\[
\log\abs{\det J}=\sum_j \left[s_\theta(x_a)\right]_j .
\]

代价从 \(O(d^3)\) 降到把 \(s_\theta\) 的输出加起来。求逆同样是闭式的——已知 \(y_a=x_a\)，就能重算 \(s_\theta,t_\theta\)，然后 \(x_b=(y_b-t_\theta)\odot\exp(-s_\theta)\)，不需要任何迭代求解。\(s_\theta\) 的复杂度没有让可逆性变难，因为它从来不需要被求逆。

这个技巧的限制写在第一行：\(y_a=x_a\)，一层耦合只动了一半变量。要让每个维度都被别的维度影响过，就必须堆很多层，并且在层与层之间交换划分方式（棋盘掩码与通道掩码交替，或插入固定置换、可逆的 \(1\times1\) 卷积）。深度不再是为了表达力富余，而是为了补上单层结构性的残缺。整个流的对数行列式是各层的和，反向传播照常。

\begin{remark}
交替划分是必需的，不是锦上添花。若所有层都用同一种划分，\(x_a\) 这半边从头到尾原样穿过整个网络，\(f_\theta\) 在这些坐标上是恒等映射，模型对它们的边缘分布只能照抄基分布。这类错误在代码里不报错，只表现为似然停在一个说不清的水平上。
\end{remark}

\subsection{精确的似然买到了什么，没买到什么}

买到的是一个可比较、有单位、可以做似然比的数，以及一个可以真正采样的模型——从 \(\mathcal N(0,I)\) 抽 \(z\)，过一遍 \(g_\theta\) 就是一个样本，没有 MCMC，没有多步迭代。同一个网络同时给出密度和采样器，这在四条路线里是独一份。

没买到的是感知质量。流站在前向 KL 的方向上，覆盖式的性格一并继承：它不敢在任何有数据的地方留空洞，于是把质量摊得很开，包括摊到数据流形附近那些什么也不是的区域。有一个被反复复现的经验观察值得记住——在 CIFAR-10 上训练的流模型，对 SVHN 的图片给出的似然\emph{更高}，尽管 SVHN 根本不在训练分布里。这是经验规律，不是定理，但它足以说明一件事：不能把高似然直接读成``模型认得这个样本''，更不能拿它当异常检测的现成判据。开篇那个想法因此并没有被完全满足——你确实拿到了一个有单位的数，但这个数衡量的不是你以为的那件事。这个落差在本单元第四篇会被单独拆开。

也没买到架构自由。为了行列式便宜，每一层的结构都被限定；为了可逆，维数一格都不能省。同样的参数量下，流的表达力通常明显不如无约束的网络。

下一条路线把这两笔账反过来结：架构完全放开，可逆性不要，密度也不写——只留一个采样器，然后想办法在没有密度的情况下把散度估出来。
